% 本轮只修订源文件, 不构建 PDF.
\documentclass[UTF8]{ctexart}
\usepackage{amsmath,amssymb,amsthm}
\usepackage[margin=2.5cm]{geometry}
\usepackage[hyphens]{url}
\usepackage[hidelinks]{hyperref}
\usepackage{booktabs}
\emergencystretch 3em
\newtheorem{theorem}{定理}[section]
\newtheorem{lemma}[theorem]{引理}
\newtheorem{corollary}[theorem]{推论}
\newtheorem{remark}[theorem]{注}
\newtheorem{definition}[theorem]{定义}
\newcommand{\eps}{\epsilon}
\newcommand{\sig}{\sigma}

\title{两组明确三阶递推的精确理论:\\阶乘降阶, 最小解与有理比值分类}
\author{研究证明文档 (项目: BVE research)}
\date{2026-09-21 作者修订稿}
\begin{document}
\maketitle

\begin{abstract}
对本文逐项列出的偶次和奇次系数, 固定 $c>0$, 阶乘缩放把三阶递推
精确化为二阶差分的一阶递推 $d_j=\theta_jd_{j-1}$.
由此得到全部解 $v_j=A+Bj+C\Phi_j$, 正项有限向后公式, 唯一归一化
最小解, 以及两个奇偶分支的一般 $c$ 渐近常数闭式.
有限极点论证证明: 相邻比值在充分大整数上为有理函数当且仅当 $C=0$.
这同时排除最小解的有理比值, 不设有理次数上限.
文中完整列出首一二次分子/分母的所有参数表示, 包括可约分支;
一般降阶公式补足逐项非零假设和反向初值.
结果仅限于所列系数, 不外推到任意三阶系数族或非齐次源项.
\end{abstract}

\noindent\textbf{证明与验收范围.}
以下为数学修补作者重新推导并写全的候选证明, 需另行新上下文独立复审后发布.
精确代数检查及有限指标测试只是辅助证据, 不代替下文的全指标证明.
本稿不宣称完整 Lean 形式化, 不改变 canonical 接收状态.
第 \ref{sec:history} 节单独保留未重新审计的历史形式级数和数值记录.

\medskip
\noindent\textbf{来源区分.}
第四轮输入 \path{research/artifacts/proof-audit-round4-20260921/submitted/constructive_repairs.md}
的 B--F 提供了待核实的阶乘/正项尾和/极点路线;
本稿据实际系数重新推导, 补全终端尺度、定义域和全部可约参数表示.
2026-08-22 的根-1 高次排除证明确已存在, 其原审计及范围见
第 \ref{sec:oldproof} 节. 历史文件保持原状, 其旧状态不充当本稿的新验收.

\tableofcontents

\section{系数、尺度与定义域}

令 $\eps=0,1$ 分别表示偶次和奇次, $c>0$ 为固定实数. 序列可取复数值.
全文研究且仅研究
\begin{equation}\label{eq:mu}
 c^2\mu_j=P_j\mu_{j-1}-Q_j\mu_{j-2}+R_j\mu_{j-3},\qquad j\ge3,
\end{equation}
其中
\begin{align}\label{eq:coeff}
 P_j&=4cj(2j+2\eps-1)+c^2\frac{j}{j-1},\notag\\
 Q_j&=4j(j-1)(2j+2\eps-1)(2j+2\eps-3)+4cj(2j+2\eps-3),\\
 R_j&=4j(j-2)(2j+2\eps-3)(2j+2\eps-5).\notag
\end{align}
前三项自由, 不附加 $j=2$ 方程. 对 $j\ge3$, $R_j>0$,
因而既能前推, 也能由连续三项唯一向后恢复到零项.

设
\begin{equation}\label{eq:scales}
 z_j=\frac{\mu_j}{(j!)^2(4/c)^j},\quad
 v_j=\frac{c^j\mu_j}{(2j+\eps)!},\quad
 H_j=\frac{(2j+\eps)!}{4^j(j!)^2},\quad
 \sig=\eps-\frac12.
\end{equation}
于是 $z_j=H_jv_j$, $H_0=1$, 且
\begin{equation}\label{eq:Hratio}
 \frac{H_j}{H_{j-1}}=1+\frac{\sig}{j}>0\quad(j\ge1).
\end{equation}
所有尺度转换可逆. $z$ 的递推为
\begin{equation}\label{eq:z}
 z_j=a_1(j)z_{j-1}+a_2(j)z_{j-2}+a_3(j)z_{j-3},\qquad j\ge3,
\end{equation}
其中
\[
 a_1=\frac{P_j}{4cj^2},\quad
 a_2=-\frac{Q_j}{16j^2(j-1)^2},\quad
 a_3=\frac{cR_j}{64j^2(j-1)^2(j-2)^2}.
\]
系数趋于 $(2,-1,0)$, 极限特征多项式为 $\rho(\rho-1)^2$.
这里不以这一极限代替向后收敛或渐近定理的证明.

\begin{lemma}[比值等价的非零前提]\label{lem:fp}
 若 $E_j\ne0$ 对所有 $j\ge0$, 则 $E$ 解 \eqref{eq:z} 当且仅当
 $e_j=E_j/E_{j-1}$ 满足
\begin{equation}\label{eq:ratio}
 e_j=F_j(e_{j-1},e_{j-2}),\qquad
 F_j(x,y)=a_1(j)+\frac{a_2(j)}x+\frac{a_3(j)}{xy},\quad j\ge3.
\end{equation}
 若另有 $E_0=1$, 则 $E_j=\prod_{k=1}^je_k$.
 对一个最终无零项的序列, 同一结论只在所用各项均非零的尾部成立.
\end{lemma}
\begin{proof}
 将 \eqref{eq:z} 除以非零的 $E_{j-1}$ 即得 \eqref{eq:ratio};
 反向乘回. 连乘比值给出乘积式.
 尾部论证要求 $E_{j-3},\ldots,E_j$ 均非零, 不对有限前缀作推断.
\end{proof}

\section{精确阶乘降阶与全部解}

\begin{theorem}[二阶差分因式分解]\label{thm:factor}
 设
\[
 \theta_j=\frac{c}{2(j-1)(2j+2\eps-1)}\quad(j\ge2),\qquad
 a_j^{\mathrm{tail}}=\frac{jc^j}{(2j+\eps)!}\quad(j\ge1).
\]
 则 \eqref{eq:mu} 等价于
\begin{equation}\label{eq:v}
 v_j=(2+\theta_j)v_{j-1}-(1+2\theta_j)v_{j-2}+\theta_jv_{j-3},\quad j\ge3.
\end{equation}
 若 $d_j=v_j-2v_{j-1}+v_{j-2}$ ($j\ge2$), 则进一步等价于
\begin{equation}\label{eq:d}
 d_j=\theta_jd_{j-1}\quad(j\ge3),\qquad
 d_j=C a_j^{\mathrm{tail}}\quad(j\ge2)
\end{equation}
 其中 $C$ 为任意常数.
\end{theorem}
\begin{proof}
 对 $\eps=0$ 或 $1$, 分别消去连续阶乘因子得到
\begin{align*}
 \frac{P_j(2j-2+\eps)!}{c(2j+\eps)!}&=2+\theta_j,\\
 \frac{Q_j(2j-4+\eps)!}{(2j+\eps)!}&=1+2\theta_j,\\
 \frac{cR_j(2j-6+\eps)!}{(2j+\eps)!}&=\theta_j.
\end{align*}
 例如 $(2j+\eps)(2j+\eps-1)=2j(2j+2\eps-1)$;
 后两式同样用此恒等式及其平移形式逐对约去因子.
 将 \eqref{eq:v} 移项就是 $d_j-\theta_jd_{j-1}=0$.
 又
\[
 \frac{a_j^{\mathrm{tail}}}{a_{j-1}^{\mathrm{tail}}}
 =\frac{cj}{(j-1)(2j+\eps)(2j+\eps-1)}=\theta_j\quad(j\ge2).
\]
 取 $C=d_2/a_2^{\mathrm{tail}}$, 归纳得到 \eqref{eq:d}; 反向显然成立.
 $\theta_2$ 的定义用于尾权重比, 不添加原递推的 $j=2$ 条件.
\end{proof}

\begin{theorem}[正项尾和与通解]\label{thm:general}
 对 $j\ge0$, 定义
\begin{equation}\label{eq:Phi}
 \Phi_j=\sum_{r=j+2}^{\infty}(r-j-1)a_r^{\mathrm{tail}}.
\end{equation}
 则 $\Phi_j>0$, $\Phi_j\to0$, 并且全部解恰为
\begin{equation}\label{eq:general}
 v_j=A+Bj+C\Phi_j,\qquad z_j=H_j(A+Bj+C\Phi_j).
\end{equation}
 三个常数 $A,B,C$ 唯一.
\end{theorem}
\begin{proof}
 因为 $a_{r+1}^{\mathrm{tail}}/a_r^{\mathrm{tail}}=O(r^{-2})$,
 $\sum_{r\ge1}r a_r^{\mathrm{tail}}$ 收敛.
 因此 \eqref{eq:Phi} 绝对收敛, 严格为正,
 且受该收敛正项级数的尾部控制而趋于零.
 相减得
\[
 \Phi_j-\Phi_{j-1}=-\sum_{r=j+1}^{\infty}a_r^{\mathrm{tail}}\quad(j\ge1),
 \qquad \Phi_j-2\Phi_{j-1}+\Phi_{j-2}=a_j^{\mathrm{tail}}\quad(j\ge2).
\]
 由 \eqref{eq:d}, $v_j-C\Phi_j$ 二阶差分为零, 故为 $A+Bj$.
 反过来该表示直接满足 \eqref{eq:v}.
 $C$ 由 $d_2$ 决定, 再由零项和一项决定 $A,B$, 因而唯一.
 等价地, $1,j,\Phi_j$ 在指标 $0,1,2$ 的值矩阵行列式为
 $\Phi_2-2\Phi_1+\Phi_0=a_2^{\mathrm{tail}}>0$.
\end{proof}

\begin{theorem}[两个显式非最小解]\label{thm:closed}
 归一化 $E_0^-=E_0^+=1$ 的两个解为
\[
 E_j^-=H_j,\qquad E_j^+=\frac{j+\sig+1}{\sig+1}H_j.
\]
 其相邻比值分别为 $e_j^-=1+\sig/j$ 和 $e_j^+=1+(\sig+1)/j$.
 在 $\mu$ 尺度下, 偶次为
\[
 \mu_j^-=(2j)!/c^j,\qquad \mu_j^+=(2j+1)!/c^j;
\]
 奇次为
\[
 \mu_j^-=(2j+1)!/c^j,\qquad \mu_j^+=(2j+3)!/(6(j+1)c^j).
\]
 它们在 $j\ge3$ 解原递推, 各项严格为正.
\end{theorem}
\begin{proof}
 在 \eqref{eq:general} 中分别取 $v=1$ 和 $v=(j+\sig+1)/(\sig+1)$.
 利用 \eqref{eq:Hratio} 计算比值, 再按 \eqref{eq:scales} 还原阶乘即得.
 这里 $\sig+1=1/2$ 或 $3/2$, 从而正性及归一化均成立.
\end{proof}

\section{有限向后公式与唯一最小解}

\begin{definition}
 非零解 $h$ 称为最小解, 若对每个与它线性无关的解 $z$,
 $h_j/z_j$ 在充分大指标有定义且趋于零.
 同一非零尺度乘到所有解上不改变此定义.
\end{definition}

\begin{theorem}[终端尺度、向后极限与最小性]\label{thm:minimal2}
 固定整数 $N\ge2$, 取终端数据
 $z_N^{(N)}=1$, $z_{N+1}^{(N)}=z_{N+2}^{(N)}=0$,
 用 \eqref{eq:z} 从 $j=N+2$ 反推到 $j=3$. 记
\[
 \Phi_j^{(N)}=\sum_{r=j+2}^{N+2}(r-j-1)a_r^{\mathrm{tail}}
 \quad(0\le j\le N+2),
\]
 空和为零. 则未归一化的精确公式为
\begin{equation}\label{eq:finite}
 v_j^{(N)}=\frac{\Phi_j^{(N)}}{H_N a_{N+2}^{\mathrm{tail}}},\qquad
 z_j^{(N)}=\frac{H_j\Phi_j^{(N)}}{H_N a_{N+2}^{\mathrm{tail}}}
 \quad(0\le j\le N+2).
\end{equation}
 特别地 $z_0^{(N)}>0$, 且对每个固定 $j\ge0$,
\begin{equation}\label{eq:minimal}
 h_j^*:=\lim_{N\to\infty}\frac{z_j^{(N)}}{z_0^{(N)}}
 =\frac{H_j\Phi_j}{\Phi_0}>0,\qquad
 \mu_j^*=\frac{(2j+\eps)!}{c^j}\frac{\Phi_j}{\Phi_0}.
\end{equation}
 这是唯一满足 $h_0^*=\mu_0^*=1$ 的最小解.
\end{theorem}
\begin{proof}
 有限和也满足 $\Phi_j^{(N)}-2\Phi_{j-1}^{(N)}+\Phi_{j-2}^{(N)}
 =a_j^{\mathrm{tail}}$ ($2\le j\le N+2$).
 因而 \eqref{eq:finite} 满足所需递推.
 $\Phi_N^{(N)}=a_{N+2}^{\mathrm{tail}}$,
 $\Phi_{N+1}^{(N)}=\Phi_{N+2}^{(N)}=0$, 所以终端恰为 $(1,0,0)$
 的 $z$ 数据, 而不是把 $\mu_N=1$ 与 $z_N=1$ 混用.
 向后唯一性保证这是指定解. 也可从
 $d_{N+2}^{(N)}=v_N^{(N)}=1/H_N$ 反推得到同一常数.

 因 $H_0=1$, 零项为 $\Phi_0^{(N)}/(H_N a_{N+2}^{\mathrm{tail}})>0$.
 归一化后常数约去, 分子分母的正项有限和各自趋于
 $\Phi_j$ 和 $\Phi_0>0$, 证明固定指标收敛.

 在通解 \eqref{eq:general} 中, $A=B=0$ 给出一维的衰减 $v$ 分支.
 其余解的 $A+Bj$ 不恒零, 因 $\Phi_j\to0$, 它们最终非零,
 且 $\Phi_j/(A+Bj+C\Phi_j)\to0$.
 因此 $H\Phi$ 为最小解. 任何含非零仿射部分的解都不能相对 $H\Phi$
 趋于零, 故不最小. 最后由 $\Phi_0>0$ 得唯一归一化.
\end{proof}

\section{两种奇偶性的一般渐近常数}\label{sec:K}

\begin{theorem}[严格尾估计与闭式]\label{thm:K}
 对固定 $c>0$, 有
\begin{align}\label{eq:Kasym}
 \mu_j^*&=K_\eps(c)j^{-3}
  \left(1-\frac{2\eps+3}{j}+O(j^{-2})\right),\notag\\
 h_j^*&=K_\eps(c)\frac{(c/4)^j}{(j!)^2}j^{-3}(1+O(j^{-1})),\\
 \frac{h_j^*}{h_{j-1}^*}&=\frac{c}{4j^2}
  \left(1-\frac3j+O(j^{-2})\right),\notag
\end{align}
 其中隐含常数可依赖 $c,\eps$, 并且
\begin{align}\label{eq:Kclosed}
 K_0(c)&=\frac{c^2}{4(c\cosh\sqrt c-\sqrt c\sinh\sqrt c)},\notag\\
 K_1(c)&=\frac{c^{5/2}}
 {4((c+3)\sinh\sqrt c-3\sqrt c\cosh\sqrt c)}.
\end{align}
 两个分母在 $c>0$ 都严格为正. 特别地
 $K_0(1)=e/4$, $\lim_{c\downarrow0}K_0(c)=3/4$,
 $\lim_{c\downarrow0}K_1(c)=15/4$.
 后两个式子是右极限, 不把 $c=0$ 纳入原递推.
\end{theorem}
\begin{proof}
 在 \eqref{eq:minimal} 的和中令 $r=j+\ell+1$. 第一项为
\[
 T_{j,1}=\frac{c^2(j+2)}{\Phi_0}
 \frac{(2j+\eps)!}{(2j+4+\eps)!}.
\]
 对 $\ell\ge2$, 第 $\ell$ 项与第一项之比不超过
\[
 \ell^2 q_j^{\ell-1},\qquad q_j=\frac{c}{(2j+5+\eps)^2}.
\]
 这是因为 $(j+\ell+1)/(j+2)\le\ell$, 新增的 $2\ell-2$ 个
 阶乘因子均至少为 $2j+5+\eps$.
 对充分大的 $j$, $q_j<1$, 且
\[
 \sum_{\ell=2}^{\infty}\ell^2q_j^{\ell-1}
 =\frac{1+q_j}{(1-q_j)^3}-1=O(j^{-2}).
\]
 故 $\mu_j^*=T_{j,1}(1+O(j^{-2}))$.
 展开第一项的四个线性因子给出
\[
 T_{j,1}=\frac{c^2}{16\Phi_0}j^{-3}
  \left(1-\frac{2\eps+3}{j}+O(j^{-2})\right).
\]
 于是 $K_\eps(c)=c^2/(16\Phi_0)$. 相邻两项相除并使用
 $z_j/z_{j-1}=(c/(4j^2))(\mu_j/\mu_{j-1})$, 得到 \eqref{eq:Kasym}.
 此估计还说明 $\Phi_j$ 比任意负幂衰减更快.

 为计算 $\Phi_0$, 令 $D=c\,d/dc$. 绝对局部一致收敛的整函数级数为
\[
 f_0(c)=\sum_{r=0}^{\infty}\frac{c^r}{(2r)!}=\cosh\sqrt c,\qquad
 f_1(c)=\sum_{r=0}^{\infty}\frac{c^r}{(2r+1)!}=\frac{\sinh\sqrt c}{\sqrt c}.
\]
 这些初等级数亦见 \cite{dlmf}. 逐项求导合法, 因而
\begin{align*}
 \Phi_0\big|_{\eps=0}=D(D-1)f_0(c)
  &=\frac{c\cosh\sqrt c-\sqrt c\sinh\sqrt c}{4},\\
 \Phi_0\big|_{\eps=1}=D(D-1)f_1(c)
  &=\frac{(c+3)\sinh\sqrt c-3\sqrt c\cosh\sqrt c}{4\sqrt c}.
\end{align*}
 两者均为正项级数 $\sum_{r\ge2}r(r-1)c^r/(2r+\eps)!>0$,
 证明分母正性及 \eqref{eq:Kclosed}.
 在偶次 $c=1$, $\Phi_0=(\cosh1-\sinh1)/4=1/(4e)$.
 当 $c\downarrow0$, $\Phi_0=2c^2/(4+\eps)!+O(c^3)$,
 从而 $K_\eps(c)\to(4+\eps)!/32$, 即所述两极限.
\end{proof}

\begin{remark}[旧锚点的范围]
 \path{docs/SL_third_order_K1_proof.tex} 已给出偶次 $c=1$ 的独立文档证明.
 该文的有限终端为 $\mu_{N+1}=1$, $\mu_N=\mu_{N-1}=0$;
 本文定理 \ref{thm:minimal2} 的有限终端不同, 已直接证明自己的归一化公式.
 两者的最终最小解一致. 本文一般 $c$、双奇偶的证明不把旧锚点的历史验收
 扩大为新增结论的验收.
\end{remark}

\section{一般非零基准下的精确降阶}

\begin{theorem}[降阶公式及反向初值]\label{thm:reduction}
 设 $E$ 解 \eqref{eq:z}, 且 $E_j\ne0$ 对每个 $j\ge0$.
 对任意解 $z$, 令 $r_j=z_j/E_j$ ($j\ge0$),
 $s_j=r_j-r_{j-1}$ ($j\ge1$). 则
\begin{equation}\label{eq:reduction}
 s_j=A_js_{j-1}+B_js_{j-2},\quad
 A_j=-\frac{a_2(j)E_{j-2}+a_3(j)E_{j-3}}{E_j},\quad
 B_j=-\frac{a_3(j)E_{j-3}}{E_j},\quad j\ge3.
\end{equation}
 反之, 任取 $r_0,s_1,s_2$, 用 \eqref{eq:reduction} 生成 $s_j$ ($j\ge3$),
 再定义 $r_j=r_0+\sum_{k=1}^js_k$, $z_j=E_jr_j$, 就得到原递推的解.
 若仅在一个尾部有逐项非零, 同样结论只在相应尾部成立, 初值指标随之平移.
\end{theorem}
\begin{proof}
 将 $z_j=E_jr_j$ 代入递推, 并以
 $r_{j-2}=r_{j-1}-s_{j-1}$,
 $r_{j-3}=r_{j-1}-s_{j-1}-s_{j-2}$ 消去各项.
 因 $E$ 自身解递推, 得到正确的中间恒等式
\begin{equation}\label{eq:intermediate}
 0=E_js_j+(a_2E_{j-2}+a_3E_{j-3})s_{j-1}+a_3E_{j-3}s_{j-2}.
\end{equation}
 除以 $E_j$ 即得 \eqref{eq:reduction}.
 反向构造给出 $z_0=E_0r_0$,
 $z_1=E_1(r_0+s_1)$, $z_2=E_2(r_0+s_1+s_2)$;
 逐步逆用 \eqref{eq:intermediate} 证明所有 $j\ge3$ 的原方程.
\end{proof}

\begin{remark}[非零解与逐项非零不同]
 偶次 $E_j=(1-j/3)H_j$ 满足 $E_0=1$, 但 $E_3=0$,
 不能用于跨过指标 $3$ 的除法. $E^\pm$ 则确实逐项为正.
 旧脚本 \path{scripts/h3_v56_odd_explicit.py} 的另一降阶式不是
 \eqref{eq:reduction}; 这里的证明不依赖该脚本.
\end{remark}

\begin{theorem}[变差常数的完整初始化]\label{thm:minimal}
 沿用定理 \ref{thm:reduction}, 设其二阶递推有已知解
 $s_j^-\ne0$ 对所有 $j\ge1$. 取 $w_2=1$, $t_1=0$, 并定义
\[
 w_j=-\frac{B_js_{j-2}^-}{s_j^-}w_{j-1}\quad(j\ge3),\quad
 t_j=\sum_{k=2}^jw_k\quad(j\ge2),\quad s_j^{\mathrm{ind}}=s_j^-t_j.
\]
 则 $s^{\mathrm{ind}}$ 与 $s^-$ 为两个独立解. 用任意 $r_0$ 和
 $r_j=r_0+\sum_{k=1}^js_k^{\mathrm{ind}}$ 得到第三阶解.
 取 $E=E^+$ 及 $s_j^-=E_j^-/E_j^+-E_{j-1}^-/E_{j-1}^+$ 时,
 可得到基 $E^+,E^-,z^{\mathrm{ind}}$.
\end{theorem}
\begin{proof}
 置 $s_j=s_j^-t_j$, 用已知解方程消去 $t_{j-1}$, 得
 $s_j^-(t_j-t_{j-1})=-B_js_{j-2}^-(t_{j-1}-t_{j-2})$.
 这正是 $w$ 方程. 两解在 $j=1,2$ 的行列式为
 $s_1^-s_2^-\ne0$, 因为 $t_1=0,t_2=1$.
 又 $B_j\ne0$, 其 Casoratian 随指标递推也不为零.
 对 $E=E^+$, 比值为 $(\sig+1)/(j+\sig+1)$, 所以
\[
 s_j^-=-\frac{\sig+1}{(j+\sig)(j+\sig+1)}\ne0\quad(j\ge1).
\]
 取比值差的线性映射 $z\mapsto s$ 的核恰为 $\operatorname{span}\{E^+\}$,
 两个独立的 $s$ 因而给出三个独立的 $z$.
\end{proof}

\section{不设次数上限的有理比值分类}

\begin{definition}
 有理比值解指一个非零解 $z$, 存在 $r\in\mathbb C(x)$,
 使充分大整数 $j$ 上 $z_{j-1}\ne0$, $r(j)$ 无极点且
 $z_j/z_{j-1}=r(j)$. 这是最终定义的比值, 不自动包含 $E_0=1$
 或从第 $1$ 项起处处有定义的乘积轨迹.
\end{definition}

\begin{theorem}[全有理比值分类, 含最小支排除]\label{thm:full}
 对 \eqref{eq:coeff} 的系数, 通解 \eqref{eq:general} 为有理比值解
 当且仅当 $C=0$ 且 $(A,B)\ne(0,0)$. 全部比值为
\begin{equation}\label{eq:rational}
 r(j)=\frac{j+\sig}{j}\frac{A+Bj}{A+B(j-1)}.
\end{equation}
 最小解的相邻比值不是有理函数. 每个有理比值均趋于 $1$,
 约分后的分子和分母次数至多为 $2$.
\end{theorem}
\begin{proof}
 非零解最终无零项: 若仿射部分非零, 它最终支配 $C\Phi_j$;
 否则它是非零常数乘以正的 $\Phi_j$.
 由于 $H_j/H_{j-1}=(j+\sig)/j$ 有理且最终非零,
 $z$ 有理比值等价于 $q(j)=v_j/v_{j-1}$ 有理.

 假设 $C\ne0$ 且 $q$ 有理. 在充分大整数上除以 $v_j$ 得
\[
 \frac{C a_j^{\mathrm{tail}}}{v_j}
 =1-\frac2{q(j)}+\frac1{q(j)q(j-1)}=:S(j).
\]
 左边不为零, 故右边不是零有理函数. 令 $R(x)=1/S(x)$,
 就有 $v_j=C a_j^{\mathrm{tail}}R(j)$ 在充分大整数上成立.
 设
\[
 D(x)=\frac{2(x-1)(2x+2\eps-1)}c,
\]
 则 $a_{j-1}^{\mathrm{tail}}/a_j^{\mathrm{tail}}=D(j)$.
 将表示代入二阶差分并除以 $Ca_j^{\mathrm{tail}}$ 得
\begin{equation}\label{eq:pole}
 R(x)-2D(x)R(x-1)+D(x)D(x-1)R(x-2)=1.
\end{equation}
 它在无穷多个整数成立, 因而为有理函数恒等式.

 若 $R$ 有有限复极点, 取实部最小的一个 $p$.
 在 $x=p$, $R(x-1)$ 和 $R(x-2)$ 均解析, 否则 $p-1$ 或 $p-2$
 是实部更小的极点. 多项式系数不能引入新极点,
 故系数为 $1$ 的 $R(x)$ 的极点无法抵消, 与右边为常数矛盾.
 因而 $R$ 是非零多项式. 若次数为 $m\ge0$, \eqref{eq:pole} 左边
 三项次数分别为 $m,m+2,m+4$; 最后一项最高次系数为
 $16/c^2$ 乘以 $R$ 的首项系数, 非零, 也不可能等于 $1$.
 矛盾说明 $C=0$.

 反之, 非零仿射 $v=A+Bj$ 解递推, 在充分大指标不为零,
 且比值正是 \eqref{eq:rational}. 所余结论由此立即得到.
\end{proof}

\subsection{归一化轨迹与最终尾部}

当 $C=0$ 时 $z_0=A$. 因而 $z_0=1$ 必须且只需 $A=1$.
$B=0$ 给出 $E^-$; $B\ne0$ 时设 $\tau=A/B-1$, 得
\begin{equation}\label{eq:tau}
 E_j^{(\tau)}=\frac{j+\tau+1}{\tau+1}H_j,\quad \tau\ne-1,
 \qquad r_\tau(j)=\frac{(j+\sig)(j+\tau+1)}{j(j+\tau)}.
\end{equation}
其 $E^+,E^-$ 系数分别为
 $(\sig+1)/(\tau+1)$ 与 $(\tau-\sig)/(\tau+1)$.
若要求从第 $1$ 项起所有相邻比值都有定义, 必须排除全部
$\tau\in\{-1,-2,-3,\ldots\}$; 其中 $\tau=-1$ 首先就不允许归一化,
其余值使某个正指标项为零. 序列本身在 $\tau=-2,-3,\ldots$ 仍是合法解,
其零点之后可以使用尾部比值.

允许最终比值时还必须包含 $A=0,B\ne0$, 即 $v=Bj$ 的 $\tau=-1$ 分支.
它的约分比值为 $(j+\sig)/(j-1)$, $z_0=0$, 不能乘常数变成 $z_0=1$.
也不存在有相同尾部比值而零项为 $1$ 的别的全局解:
相同比值的非零尾部成比例, 向后唯一性把比例延伸到零项.
有理函数约分后的可去奇点与原序列的零项、或未经约分的表达式的零分母,
是三个不同问题; 约分不授权在原式零分母处直接作除法.

\subsection{全部首一二次参数表示}

\begin{theorem}[四参数表, 包括所有约分]\label{thm:free}
 写
\[
 r(x)=\frac{x^2+ax+b}{x^2+\gamma x+d},\qquad a,b,\gamma,d\in\mathbb C.
\]
 它在最终定义意义下满足 \eqref{eq:ratio} 当且仅当四元组属于下表.
 $t$ 为任意复数, $\tau$ 只在第一行排除两个所列值.
\begin{center}
\begin{tabular}{c|cccc}
约分后的比值 & $a$ & $b$ & $\gamma$ & $d$\\\hline
$r_\tau$, $\tau\notin\{\sig,-1\}$
 & $\sig+\tau+1$ & $\sig(\tau+1)$ & $\tau$ & $0$\\
$e^-=(x+\sig)/x$ & $\sig+t$ & $\sig t$ & $t$ & $0$\\
$e^+=(x+\sig+1)/x$ & $\sig+1+t$ & $(\sig+1)t$ & $t$ & $0$\\
$r_{-1}=(x+\sig)/(x-1)$ & $\sig+t$ & $\sig t$ & $t-1$ & $-t$
\end{tabular}
\end{center}
 若只讨论实参数, 将 $t,\tau$ 限为实数即可.
 对原二次表达式逐项求值时仍要求 $j^2+\gamma j+d\ne0$;
 表格本身分类的是有理恒等式及其全部参数表示.
\end{theorem}
\begin{proof}
 一个满足有理恒等式的非零有理函数趋于 $1$, 故可在足够大的整数处
 避开其零点和极点, 连乘构成尾部解, 再由 $a_3\ne0$ 向后延拓.
 定理 \ref{thm:full} 因而适用.
 当 $B=0$, 约分形式为 $e^-$.
 当 $B\ne0$, 为 $r_\tau$ ($\tau=A/B-1$, 可等于 $-1$).
 分子因子为 $x+\sig,x+\tau+1$, 分母因子为 $x,x+\tau$.
 因 $\sig=\pm1/2\ne0$, 唯一可能的约分是
 $\tau=\sig$ 或 $\tau=-1$, 分别给出 $e^+$ 和 $r_{-1}$.
 这两值不同, 且所得一次分子分母互素.

 若约分形式 $N_0/D_0$ 是互素首一二次式, 另一首一二次表示
 $N/D$ 满足 $ND_0=N_0D$, 互素性迫使 $N=N_0,D=D_0$.
 若约分形式是互素首一次式, 则必为
 $N=N_0(x+t),D=D_0(x+t)$, 其中 $t$ 任意.
 逐项展开这三种一次形式即为表格后三行. 这证明穷尽性及反向成立.
\end{proof}

\begin{remark}[R4-F03 的具体边界]
 偶次 $e^+$ 行是 $(a,b,\gamma,d)=(t+1/2,t/2,t,0)$,
 不是带绝对值的分支; $t=-3/4$ 给出审计中的
 $(-1/4,-3/8,-3/4,0)$.
 奇次 $e^-$ 行是 $(t+1/2,t/2,t,0)$, 符号为正.
 奇次 $e^+$ 行则是 $(t+3/2,3t/2,t,0)$, 不与奇次 $e^-$ 混用.
 最后一行也不能省略: 偶次取 $t=2$ 给出
 $(3/2,-1,1,-2)$, 约分后为 $(x-1/2)/(x-1)$;
 它是合法尾部比值, 尽管 $d=-2\ne0$ 且不能归一化为 $E_0=1$.
\end{remark}

\subsection{乘积族与渐近结论的正确范围}

\begin{theorem}[单参数乘积族]\label{thm:beta}
 固定 $c>0$, 令 $E_0(\beta)=1$,
 $E_j(\beta)=\prod_{k=1}^j(1+\beta/(2k))$.
 它解 \eqref{eq:z} 当且仅当
 $\beta\in\{2\eps-1,2\eps+1\}$.
 对这两个值, 所有因子非零, 比值等价也处处成立.
\end{theorem}
\begin{proof}
 若 $\beta=-2m$ ($m\ge1$ 为整数), 则从指标 $m$ 起全为零;
 三个连续尾部零项借助 $a_3\ne0$ 反推迫使 $E_0=0$, 矛盾.
 其余 $\beta$ 的所有因子均非零, 故由定理 \ref{thm:full} 分类.
 $B=0$ 给出 $\beta/2=\sig$.
 $B\ne0$ 时, 将 $1+\beta/(2x)$ 与 $r_\tau(x)$ 比较, 得
\[
 (x+\beta/2)(x+\tau)=(x+\sig)(x+\tau+1).
\]
 一次项给 $\beta/2=\sig+1$, 常数项给 $\tau=\sig$.
 两个比值由定理 \ref{thm:closed} 实现, 连乘归一化完成反向证明.
\end{proof}

\begin{corollary}[移位的一次比值]
 对 $r(x)=1+\alpha/(x+\gamma)$, 全部最终有理轨迹的参数为
\[
 (\alpha,\gamma)=(\sig,0),\quad(\sig+1,0),\quad(\sig+1,-1).
\]
 最后一支仅为 $A=0$ 的尾部, 从 $k=1$ 开始的乘积在第一因子就无定义.
\end{corollary}
\begin{proof}
 $\alpha=0$ 的常数比值 $1$ 不在定理 \ref{thm:full} 的列表中.
 对 $\alpha\ne0$, 一次分子分母互素, 定理 \ref{thm:free} 的后三种
 约分形式给出并穷尽这三对参数.
\end{proof}

\begin{theorem}[根-1 渐近无需形式级数假设]\label{thm:asym}
 任意非零解最终有相邻比值. 若 $A=B=0$, 它是根-0 最小支.
 其余都是根-1 解, 且
\[
 \lim_{j\to\infty}j\left(\frac{z_j}{z_{j-1}}-1\right)
 =\begin{cases}\sig,&B=0,\ A\ne0,\\\sig+1,&B\ne0.\end{cases}
\]
 所以偶次指数为 $-1/2,1/2$, 奇次为 $1/2,3/2$.
\end{theorem}
\begin{proof}
 最小支用 \eqref{eq:Kasym}. 对非零仿射部分,
 $\Phi_j$ 比任意负幂衰减更快, 从而
 $v_j/v_{j-1}$ 与 $(A+Bj)/(A+B(j-1))$ 的差也如此.
 $B=0$ 时后者为 $1$; $B\ne0$ 时为 $1+j^{-1}+O(j^{-2})$.
 再乘以 \eqref{eq:Hratio} 即得, 不从一个一阶极限假设完整渐近展开.
\end{proof}

\begin{corollary}[有理意义下的刚性]\label{thm:rigid}
 若根-1 解的比值有理且渐近指数为 $\sig$, 则它是常数乘以 $E^-$.
 若不要求有理, 则 $H_j(A+C\Phi_j)$ ($A\ne0,C\ne0$) 有同一个
 渐近指数和同一个全代数阶比值展开, 却不是常数乘以 $E^-$.
\end{corollary}
\begin{proof}
 有理情形由定理 \ref{thm:full} 使 $C=0$, 再用定理 \ref{thm:asym}
 使 $B=0$. 反例族的差异比任意负幂小, 而 $\Phi$ 非恒定,
 所以其比值不能与 $e^-$ 恒同; 定理 \ref{thm:full} 也直接排除其有理性.
\end{proof}

\subsection{已存在的根-1 高次证明及原审计}\label{sec:oldproof}

历史证明为
\path{runs/plugin-perf-eval/R-20260822T000000Z-a6-reuse/candidate_proof.md},
原审计为
\path{runs/plugin-perf-eval/R-20260822T000000Z-a6-audit/audit_report.md}.
它们处理根-1 有理比值的所有次数, 没有处理根-0.
原审计记录为 \texttt{REPAIRABLE\_GAP}, 指出自由支的完整三阶系数
消失及基线推导详略问题; 当前候选文件已明确写出自由参数的实现.
因此旧正文的 ``全仓库尚无高次排除证明'' 表述已过时,
但也不能把历史记录说成本轮新增最小支证明的独立通过.

该旧路线的全阶机制可直接说明如下. 令 $t=1/j$,
$\mathcal E(t)=1+ut+\sum_{m\ge2}x_mt^m$, $\mathcal A_i(t)=a_i(1/t)$,
$\mathcal E_1=\mathcal E(t/(1-t))$,
$\mathcal E_2=\mathcal E(t/(1-2t))$.
比值恒等式的乘积残差为
\[
 G=\mathcal E\mathcal E_1\mathcal E_2
   -\mathcal A_1\mathcal E_1\mathcal E_2-\mathcal A_2\mathcal E_2-\mathcal A_3.
\]
因 $\mathcal A_1=2+2\sig t+O(t^2)$,
$\mathcal A_2=-1-2\sig t+O(t^2)$,
在 $[t^{m+1}]G$ 中 $x_{m+1}$ 的系数为 $3-4+1=0$.
对 $m\ge2$, $x_m$ 只线性出现, 因为 $2m>m+1$;
其系数依次来自三个部分:
\[
 (6u+3m)+(-6m-4u-4\sig)+(2m+2\sig)=2u-m-2\sig.
\]
在 $u=\sig$ 支该系数为 $-m$, 在 $u=\sig+1$ 支为 $2-m$.
前一支由已知 $e^-$ 的存在和逐阶唯一性给出唯一展开.
后一支从 $m=3$ 起唯一, 而
\[
 r_\tau(1/t)=1+(\sig+1)t+(\sig-\tau)t^2+O(t^3)
\]
实现每个 $x_2\in\mathbb C$ (包括只作为尾部的 $\tau=-1$).
这也证明 $m=2$ 时整个方程可满足, 不只是其对角系数为零.
有理函数在无穷处有收敛 Laurent 展开, 同一展开决定同一有理函数,
所以高次有理比值已被排除. 此处的允许 $u$ 由定理 \ref{thm:asym}
直接保证; 主分类另有定理 \ref{thm:full} 的不依赖此展开的极点证明.

\section{撤回过强的盒式归纳结论}\label{sec:box}

定理 \ref{thm:full} 只分类有理比值, 不能推出所有齐次根-1 解都在
二维 $E^\pm$ 空间内. 原递推本身是齐次的, $C\Phi$ 也是其齐次解分量.
序列的线性叠加通常不使比值偏差线性叠加.
对确在 $C=0$ 的 $r_\tau$ 支, 在分母非零时有
\[
 \delta_j^{(\tau)}:=r_\tau(j)-e_j^+
 =\frac{\sig-\tau}{j(j+\tau)}.
\]
这说明该支上一处 $\delta_j=0$ 强制 $\tau=\sig$;
$\tau=\sig$ 给 $E^+$, 而 $\tau\to\infty$ 的归一化极限是 $E^-$.
这不是全部齐次根-1 轨迹的结论.

具体地, 取偶次 $c=1$, 初值 $(z_0,z_1,z_2)=(1,2,5/2)$.
有 $\rho_1=2$, $\rho_2=5/4=e_2^+$; 因 $a_3(3)=1/64$,
\[
 \rho_3-e_3^+=\frac{a_3(3)}{5/4}
       \left(\frac12-\frac{1}{3/2}\right)=-\frac1{480}.
\]
故 $\delta_1=1/2>0$, $\delta_2=0$, $\delta_3<0$.
这确实是根-1 解: $v_0=1,v_1=4,v_2=20/3$,
所以 $C=(v_2-2v_1+v_0)/a_2^{\mathrm{tail}}=-4$;
$B=3+4(\Phi_1-\Phi_0)>2$, 因 $\Phi_1>0$, $\Phi_0=1/(4e)<1/4$.
定理 \ref{thm:asym} 适用.

因此撤回原文对全部齐次根-1 轨迹的退化配置排除及由此宣称的盒封闭.
旧随机盒测试只能作为其实际采样域的历史证据, 不证明全量词前向不变性.
这也不等于否定一切加有其他条件的盒; 这样的新定理需要单独准确陈述和证明.
非齐次源项的控制同样未在这里解决.

\section{保留的历史形式级数和数值记录}\label{sec:history}

\noindent\textbf{本节未重新审计.}
以下来自旧 \path{scripts/op13_*.py} 记录, 本轮没有重跑这些程序,
也不继承其旧的 ``全符号验证'' 或精度标签.
高阶系数列在这里便于将来核对; 已证明的渐近阶数仅如定理 \ref{thm:K}.
定义一个历史截断表达式, 不声称它与实际 $\rho_j^*$ 之差具有指定余项:
\begin{equation}\label{eq:rhos}
 \widehat\rho_j=\frac{c}{4j^2}\left(1-\frac3j+\frac{C_2}{j^2}
 +\frac{C_3}{j^3}+\frac{C_4}{j^4}+\frac{C_5}{j^5}+\frac{C_6}{j^6}\right).
\end{equation}
历史偶次系数:
\[
 C_2=6,\quad C_3=-c-\tfrac{21}2,\quad C_4=\tfrac{33c}4+\tfrac{69}4,
\]
\[
 C_5=-\tfrac{c^2}4-\tfrac{163c}4-\tfrac{219}8,\qquad
 C_6=\tfrac{85c^2}{16}+\tfrac{2529c}{16}+\tfrac{681}{16}.
\]
历史奇次系数:
\[
 C_2=8,\quad C_3=-c-\tfrac{41}2,\quad C_4=\tfrac{39c}4+\tfrac{207}4,
\]
\[
 C_5=-\tfrac{c^2}4-\tfrac{241c}4-\tfrac{1039}8,\qquad
 C_6=\tfrac{95c^2}{16}+\tfrac{4843c}{16}+\tfrac{5203}{16}.
\]
来源为 \path{scripts/op13_matched_asymp3.py}; 此处不把有限逐阶匹配
提升为全阶渐近余项证明.

历史偶次常数表 (本表的小数及提取算法未重新审计):
\begin{center}\begin{tabular}{c|rrrrrrrr}
$c$ & 0.01 & 0.1 & 0.5 & 1 & 2 & 5 & 10 & 100\\\hline
$K_0(c)$ & 0.74925 & 0.74255 & 0.71367 & 0.67957 & 0.61737
 & 0.46932 & 0.30846 & 0.00252
\end{tabular}\end{center}
来源 \path{scripts/op13_K1_definitive.py}, \path{scripts/op13_K_grid.py}.
旧小 $c$ 数值拟合为
$\log K_0(c)\approx\log(3/4)-c/10+0.0014286c^2-0.0000422c^3$,
来源 \path{scripts/op13_K_cubic.py}; 此处仅保存拟合记录, 不赋予余项阶数.
一般 $c$ 的准确常数现在由 \eqref{eq:Kclosed} 给出,
不再列为本文所研究两组系数的开放问题.

旧 \path{scripts/d4_third_order_theory.py} 的闭式验证曾使用整除,
第四轮输入据此指出 R4-F02. 本稿不再声称那些旧测试全部通过;
程序修复及其实际执行由相应程序作者记录. 本文证明不依赖旧整除输出,
也不以低次数有理拟合失败证明非有理性.

\section{本稿交付边界}

\noindent\textbf{本作者实际检查.}
项目外 \path{check_recurrence.py} 检查了双奇偶的符号系数变换和四行参数族,
50 个有限向后终端问题 (550 个原始项逐项比较),
一般解的降阶及 $r_0,s_1,s_2$ 反向初始化, 变差常数,
零项/约分边界和旧错误式的负对照, 两个常数分母及端点极限.
这些共 2804 项精确或符号断言通过; 另有 32 项高精度辅助断言通过.
源文件还作环境、括号、标签和引用的静态检查, 未作 TeX 编译.
以上是本作者的新检查, 原始日志及版本绑定见项目外 \path{results.json}.
协调者另外通报的活动 d4 运行、行为测试和软件复审是独立的一组执行证据;
它们由协调者的实际日志与报告绑定, 不计入上述作者检查,
也不代替新解析证明的独立复审.

\begin{itemize}
 \item 本稿提供所列系数的完整作者推导, 包括最小支非有理性,
 一般 $c>0$ 的双奇偶常数以及全部首一二次参数表示;
 新解析证明仍需另一个新上下文审稿人验收.
 \item 第 \ref{sec:history} 节高阶形式系数、数值精度与历史程序输出未重新审计.
 不从本轮精确局部检查推断整个旧程序库或全仓库数学正确.
 \item 一般变系数族、非齐次源项与任何额外盒不变性命题均不在本稿证明内.
 与 $H^3$ 完备性、M3/KP 及其他研究主链的关系需分别核对,
 不由本文自动升级其状态.
 \item 本轮不构建 PDF, 不声称完整 Lean 证明, 不改历史工件或 canonical.
 作者检查脚本、输入哈希与日志保存在项目外授权的
 \path{/mnt/f/tools/math-audit-round4-20260921/recurrence-author/}.
\end{itemize}

\begin{thebibliography}{4}
\bibitem{round4} 第四轮审计输入, \emph{proof audit round4 20260920}
 及 \emph{constructive repairs},
 \path{research/artifacts/proof-audit-round4-20260921/submitted/}.
 作为修补线索, 不作为已验收结论.
\bibitem{root1} 根-1 高次排除候选证明及原审计,
 \path{runs/plugin-perf-eval/R-20260822T000000Z-a6-reuse/candidate_proof.md};
 \path{runs/plugin-perf-eval/R-20260822T000000Z-a6-audit/audit_report.md}.
\bibitem{anchor} 项目文档, \emph{A Strict Proof of the Even Minimal-Solution
 Anchor $K(1)=e/4$}, \path{docs/SL_third_order_K1_proof.tex}.
\bibitem{dlmf} NIST Digital Library of Mathematical Functions,
 \S4.33, equations 4.33.1--4.33.2,
 \url{https://dlmf.nist.gov/4.33} (本轮核对初等双曲函数级数).
\end{thebibliography}
\end{document}
